\documentclass[../../notes.tex]{subfiles}
\begin{document}
\chapterfigs{01-planar-motion}

\chapter{Coordinate Frames and Planar Rigid-Body Motion}\label{ch:pl}

This chapter builds the language of robot pose in the plane: frames and relative poses, planar rotations, homogeneous transformations, pose graphs, and finally the matrix exponential and twists, which describe how a pose \emph{moves}. Everything here is the two-dimensional special case of chapter~\ref{ch:rot}. Because the group is small enough to write out by hand, every idea can be seen in full.

\input{\subfix{sections/01-coordinate-frames/section}}
\input{\subfix{sections/02-rotation-2d/section}}
\input{\subfix{sections/03-pose-homogeneous/section}}
\input{\subfix{sections/04-pose-graphs/section}}
\input{\subfix{sections/05-exponentials-twists/section}}

\end{document}
